% =====================================================================
% PARTE III — Decisão metodológica e contribuição  [~12 min]
% Fonte: Chapters/estrategias_fusao, Chapters/sistema_multi_agente
% Núcleo da apresentação: é aqui que a banca decide se a tese se sustenta.
% =====================================================================

\section{Proposta}

\begin{frame}{Da matriz de decisão à arquitetura}
  \begin{itemize}
    \item Critérios de seleção metodológica.
      \pendente{matriz de decisão do Cap. 7}
    \item Por que multi-agente, e não um modelo monolítico.
  \end{itemize}
\end{frame}

% A arquitetura merece o slide inteiro: é a figura que a banca vai olhar
% mais tempo. O TikZ da tese já existe e exporta em vetorial.
\begin{frame}[plain]
  \centering
  \pendente{arquitetura\_multiagente.tex --- compilar standalone e incluir o PDF}
  % \includegraphics[width=0.92\paperwidth,height=0.84\paperheight,keepaspectratio]%
  %   {figuras/arquitetura_multiagente.pdf}
\end{frame}

\begin{frame}{Estratégias de fusão}
  \duascolunas{%
    \textbf{\color{peedark}Nível de score}
    \begin{itemize}\footnotesize
      \item \pendente{}
    \end{itemize}
  }{%
    \textbf{\color{peedark}Nível de decisão}
    \begin{itemize}\footnotesize
      \item \pendente{}
    \end{itemize}
  }
\end{frame}

\begin{frame}{Coordenação: intra-grupo, hierárquica e blackboard}
  \begin{itemize}
    \item Coordenação intra-grupo entre agentes da mesma tarefa.
    \item Coordenação hierárquica entre os três grupos.
    \item \dest{Blackboard} e meta-coordenação.
    \item \pendente{Cap. 8}
  \end{itemize}
\end{frame}

\begin{frame}{Hipóteses experimentais}
  \begin{enumerate}
    \item \pendente{H1}
    \item \pendente{H2}
    \item \pendente{H3 --- a hipótese de sinergia, verificada por ablação}
  \end{enumerate}
\end{frame}
